\chapter{狭义相对论：时间延缓}\label{ch:g12:special-relativity}

每颗 GPS 卫星都携带一台原子钟，其精度可达
三百万年仅差一秒——但发射前，工程师会有意将其调慢：若不加调整，
它每天会快 38 微秒，到明天傍晚，
你的手机就会把你定位在离你所在之处十一千米外。
错不在时钟——时间本身对运动者与静止者的流逝方式各不相同：
本章将沿着爱因斯坦的思路，从一个顽固的事实走到这一结论。

\section{有一个速率不肯相加}

速率依赖于参考系，并按加法合成
（\cref{ch:g10:relative-motion}）：一名乘客以 \qty{4}{km/h} 的速率朝车头走去，
而火车以 \qty{300}{km/h} 的速率行驶，路边观察者测得该乘客的速率为 \qty{304}{km/h}。

\begin{definition}[惯性参考系]\label{def:g12:special-relativity:inertial}
称\emph{惯性参考系}\index{惯性参考系}为一种参考系
（\cref{ch:g10:relative-motion}），在其中不受力的物体
速度保持不变：地面近似满足此条件，以及相对地面
做匀速直线运动的每一参考系。
\end{definition}

把这一规则推到极端：火车车头灯的光束应当
以 $c + \qty{300}{km/h}$ 的速率掠过地面观察者。1887 年，迈克耳孙与莫雷
比较了光沿地球 \qty{30}{km/s} 公转方向与垂直该方向的光速，
以寻找这种差异；以极高的精度，他们
却一无所获——无论哪个方向、哪个季节，光速都是同一个 $c$。爱因斯坦于 1905 年
相信了这一实验，并以之为依据重建时间与空间，基础是两条公设，
其判据在于它们的推论——一个世纪以来经受住了检验，从未失败。

\begin{theorem}[狭义相对论的两条公设]\label{thm:g12:special-relativity:postulates}
\index{狭义相对论}\index{相对论公设}
\begin{enumerate}
  \item \emph{相对性原理}：物理定律在每个惯性参考系中
        形式相同；没有哪个参考系可称“真正静止”。
  \item \emph{$c$ 的不变性}：真空中的光在每个惯性参考系中
        以 $c = \qty{3.00e8}{m/s}$ 传播，
        无论光源或观察者如何运动。
\end{enumerate}
\end{theorem}
\admitted

\section{告别“同一时刻”}

\begin{proposition}[同时性的相对性]\label{prop:g12:special-relativity:simultaneity}
\index{同时性}
在一个惯性参考系中同时发生的两个事件，一般而言
在另一惯性参考系中则不再同时：“处处此刻”依参考系而定。
\end{proposition}

\begin{proof}
爱因斯坦的火车。两道闪电分别击中运动中火车的两端，
使铁轨烧焦。站台上的观察者 $M$ 位于两处标记的正中间，
同时收到两道闪光；两道闪光各自以相同速率 $c$ 走过相同距离，
因此两次击中同时发生。乘客 $M'$ 位于车厢正中，
他\emph{迎着}车头的闪光运动，该闪光先到达他；
但在他的参考系中，两道闪光同样走过相等的距离（半列车长），
并以相同速率 $c$ 传播——这正是第二公设！——因此车头的击中先发生。
两种判断在各自的参考系中都是正确的，
这一分歧（对一列火车而言仅为纳秒量级）正是人们从未察觉的原因。
\end{proof}

\begin{omfigure}
\begin{tikzpicture}[scale=0.82, font=\small]
  \draw[thick] (0,0) -- (10.6,0);
  \fill (5.3,0) circle (2pt); \node[below] at (5.3,-0.05) {$M$};
  \node[below] at (1.6,-0.05) {$A$}; \node[below] at (9.0,-0.05) {$B$};
  \draw[omDef, thick] (1.6,0.55) rectangle (9.0,1.45);
  \fill[omDef] (5.3,1.45) circle (2pt); \node[above] at (5.3,1.5) {$M'$};
  \draw[-Stealth, omDef, thick] (9.2,1.0) -- (10.4,1.0) node[above left] {$\vect v$};
  \draw[omThm, thick] (1.25,2.7) -- (1.7,2.15) -- (1.45,2.05) -- (1.75,1.5);
  \draw[omThm, thick] (8.65,2.7) -- (9.1,2.15) -- (8.85,2.05) -- (9.15,1.5);
  \draw[dashed, omExo] (2.35,0) arc (0:60:0.75) (3.1,0) arc (0:60:1.5);
  \draw[dashed, omExo] (8.25,0) arc (180:120:0.75) (7.5,0) arc (180:120:1.5);
\end{tikzpicture}

{\small 爱因斯坦的火车：来自 $A$ 与 $B$ 的闪光同时到达 $M$；$M'$
迎着 $B$ 处的闪光行驶，先与之相遇——并断定它先击中。}
\end{omfigure}

\section{光钟与 $\gamma$ 因子}

\begin{definition}[事件，固有时]\label{def:g12:special-relativity:proper-time}
称\emph{事件}\index{事件}为在某一地点、某一时刻发生的
一件事——一次滴答、一次闪光、一次衰变。同一时钟在
两个事件处均存在时，它在二者之间所读出的时间间隔即为
\emph{固有时}\index{固有时} $\Delta t_0$。若时钟相对某一参考系运动，
则该参考系测得的时间间隔一般不同，记为 $\Delta t$。
\end{definition}

\begin{theorem}[时间延缓]\label{thm:g12:special-relativity:dilation}
\index{时间延缓}\index{$\gamma$ 因子}
若时钟以恒定速率 $v$ 穿过某一惯性参考系运动，
则在该参考系中测得其走时变慢：对于同一时钟上的两个事件，
\[
\Delta t = \gamma\, \Delta t_0, \qquad
\gamma = \frac{1}{\sqrt{1 - v^2/c^2}} \geq 1 .
\]
\end{theorem}

\begin{proof}
构造尽可能简单的时钟：两面相对的镜子相距 $L$，
两镜之间有一光子来回反射；往返一次即为一个滴答，静止时
一个滴答持续 $\Delta t_0 = 2L/c$。现将该时钟固定在一艘
以速率 $v$ 运动的船上，镜面与运动方向垂直，并从地面计量一个
滴答：光子向上行进时，镜子向侧面移动，
因此光子沿斜线飞行——同时进行两种运动，与
\cref{ch:g12:kinematics-2d}中的抛体相似。若一个滴答持续 $\Delta t$，
则每半个滴答，光子水平前进 $v\,\Delta t/2$、竖直前进 $L$；而且——
关键一步——第二公设确定了光子在地面的速率为
$c$，因此斜向半程的路径长为 $c\,\Delta t/2$。据毕达哥拉斯定理：
\[
\Bigl(\frac{c\,\Delta t}{2}\Bigr)^{\!2} = L^2 + \Bigl(\frac{v\,\Delta t}{2}\Bigr)^{\!2}
\;\Longrightarrow\;
\Delta t = \frac{2L/c}{\sqrt{1 - v^2/c^2}} = \gamma\,\Delta t_0 .
\]
与之并行运动的任何时钟——手表、心跳、衰变中的原子核——
都必须保持同步，否则这一比较就会暴露船的运动，
从而违反第一公设。
\end{proof}

\begin{omfigure}
\begin{tikzpicture}[scale=0.9, font=\small]
  \draw[thick] (-0.5,0) -- (0.5,0); \draw[thick] (-0.5,2) -- (0.5,2);
  \draw[omDef, thick, -Stealth] (0,0.05) -- (0,1.95);
  \draw[Stealth-Stealth] (0.85,0) -- (0.85,2) node[midway, right] {$L$};
  \node at (0.3,-0.75) {静止：一次滴答 $\Delta t_0 = 2L/c$};
\end{tikzpicture}
\qquad
\begin{tikzpicture}[scale=0.9, font=\small]
  \foreach \x in {0,1.4,2.8} { \draw[gray] (\x-0.4,0) -- (\x+0.4,0);
    \draw[gray] (\x-0.4,2) -- (\x+0.4,2); }
  \draw[omDef, thick, -Stealth] (0,0.03) -- (1.4,1.97);
  \draw[omDef, thick, -Stealth] (1.4,1.97) -- (2.8,0.03);
  \draw[dashed, omExo] (1.4,0) -- (1.4,2);
  \draw[-Stealth, thick] (3.5,1.0) -- (4.4,1.0) node[above left] {$\vect v$};
  \node[above left] at (0.9,0.9) {$\frac{c\,\Delta t}{2}$};
  \node[right] at (1.4,0.75) {$L$};
  \draw[Stealth-Stealth] (0,-0.35) -- (1.4,-0.35) node[midway, below] {$v\,\Delta t/2$};
  \node at (1.5,-1.15) {运动时：光子沿斜边飞行};
\end{tikzpicture}

{\small 光钟。静止时，光子上升 $L$；从地面看，
光子以同样的 $c$ 沿斜边飞行——因此滴答必须更长。}
\end{omfigure}

\begin{example}[伽马的数值]\label{ex:g12:special-relativity:gammatable}
\begin{center}\small
\begin{tabular}{lccccccc}
$v/c$    & 0.01     & 0.1   & 0.5   & 0.8   & 0.9   & 0.99 & 0.999\\
$\gamma$ & 1.00005 & 1.005 & 1.155 & 1.667 & 2.294 & 7.09 & 22.4
\end{tabular}
\end{center}
几页之内毫无动静，随后一切急剧改变：$\gamma$ 紧贴 $1$，直到
光速的三分之一；随后越过 $2$（约在 $0.9c$ 处），并在 $v \to c$ 时急剧增大。
\end{example}

\begin{omfigure}
\begin{tikzpicture}
\begin{axis}[omaxis, width=8.6cm, height=5.4cm,
    xlabel={$v/c$}, ylabel={$\gamma$},
    xmin=0, xmax=1.08, ymin=0, ymax=8,
    xtick={0,0.2,0.4,0.6,0.8,1}, ytick={1,2,...,7}]
  \addplot[omDef, thick, domain=0:0.992, samples=200] {1/sqrt(1-x^2)};
  \addplot[omExo, dashed, thick] coordinates {(1,0) (1,8)};
\end{axis}
\end{tikzpicture}

{\small $\gamma = 1/\sqrt{1 - v^2/c^2}$：在日常速率下与 $1$ 无法区分，
在渐近线 $v = c$ 处发散——这是任何有质量的物体都无法达到的速度。}
\end{omfigure}

\begin{remark}[日常生活为何从未察觉]\label{rem:g12:special-relativity:everyday}
当 $v \ll c$ 时，记 $x = v^2/c^2$：由 $(1-x)(1+x) \approx 1$ 且
$(1 + x/2)^2 \approx 1 + x$，可得 $\gamma \approx 1 + v^2/(2c^2)$。一辆汽车
以 \qty{130}{km/h} 行驶时，$\gamma - 1 \approx \num{7e-15}$（约
每行驶四百万年损失一秒）；喷气式飞机为 $\num{3.5e-13}$，
即在空中飞行九万年才损失一秒。我们的时钟当时实在太粗糙了。
\end{remark}

\section{实验的判决}

\begin{example}[μ子抵达地面]\label{ex:g12:special-relativity:muons}
宇宙射线撞击高层大气时，会产生 \emph{μ子}——
寿命很短的粒子，固有寿命 $\Delta t_0 = \qty{2.2}{\micro s}$
——在约 \qty{15}{km} 高空，速度几乎为 $c$。若不计相对论，μ子
每个寿命内只能飞行 $c\,\Delta t_0 \approx \qty{660}{m}$，下降过程
约需 $23$ 个寿命，而存活比例
$\mathrm{e}^{-23} \approx \num{e-10}$ 应使海平面的 μ子
极为罕见。然而探测器测得，约为每平方厘米
每分钟一个。时间延缓解释了这一矛盾：到达时 $\gamma \approx 30$，
在我们的时钟上，μ子存活 $30 \times \qty{2.2}{\micro s} = \qty{66}{\micro s}$，
射程 $\approx \qty{20}{km}$，足以贯穿整个大气层。
\end{example}

\begin{omfigure}
\begin{tikzpicture}[scale=0.85, font=\small]
  \draw[thick] (0,0) -- (8.6,0) node[below left] {地面};
  \draw[dashed, omExo] (0,5) -- (8.6,5);
  \node[above] at (2.1,5.02) {μ子产生于 \qty{15}{km} 高处};
  \draw[-Stealth, thick] (1.1,6.1) node[above right] {宇宙射线} -- (1.55,5.05);
  \draw[omThm, very thick, -Stealth] (3.4,5) -- (3.4,4.78);
  \node[right, omThm] at (3.5,4.62)
    {无时间延缓：$c\,\Delta t_0 \approx \qty{660}{m}$};
  \draw[omDef, very thick, -Stealth] (6.6,5) -- (6.6,0.06);
  \node[left, omDef, align=right] at (6.5,2.2)
    {$\gamma \approx 30$:\\ $\approx \qty{20}{km}$};
\end{tikzpicture}

{\small μ子的旅程：一个经典寿命在出生高度下方 \qty{660}{m} 处耗尽；
而延缓后的寿命则跨越整个大气层。}
\end{omfigure}

多项证据随之而来。1971 年，哈费尔与基廷乘客机环绕地球飞行，
携带铯原子钟（见\cref{ch:g12:oscillators-and-time}）；
降落后，它们与地面时钟的读数差为预言中的
微秒的几分之一。加速器每天都依赖时间延缓，其中的粒子
被加速到 $\gamma$ 达数千，得以存活并完成
经典意义下不可能完成的圈数。而 GPS 本身就是一个持续运行、需付出代价的实验。

\begin{remark}[GPS 的相对论账单]\label{rem:g12:special-relativity:gps}
一颗 GPS 卫星以约 \qty{3.9}{km/s} 的速率绕行，因此狭义相对论
使其时钟每天变慢约 \qty{7}{\micro s}。引力在高空较弱，
作用相反，并占了上风：广义相对论（爱因斯坦的引力理论，
在大学分册中讲述）使其时钟每天变快约
\qty{45}{\micro s}；合计，时钟每天快 \qty{38}{\micro s}。
定位是将时钟读数以速率 $c$ 换算为距离：
若不加修正，一天的误差会使位置偏差达
$c \times \qty{38}{\micro s} \approx \qty{11}{km}$——这正是开篇提到的
预先调慢时钟的原因。
\end{remark}

\begin{proposition}[长度收缩]\label{prop:g12:special-relativity:contraction}
\index{长度收缩}
在静止参考系中长为 $L_0$ 的物体，
在该物体以速率 $v$ 运动的参考系中，其沿运动方向测得的长度变短：
$L = L_0/\gamma$。
\end{proposition}
\admitted

\begin{remark}[μ子自身的视角]\label{rem:g12:special-relativity:muonframe}
长度收缩是时间延缓的另一面（严格推导见
大学分册）：μ子的钟在其自身参考系中正常计时，而迎面涌来的大气层
被压缩为 $\qty{15}{km}/30 =
\qty{500}{m}$，远在 \qty{2.2}{\micro s} 之内即穿过——两个参考系
都得出 μ子到达地面的结论。
\end{remark}

\section{质量、能量与经典世界的回归}

\begin{theorem}[质量–能量等价]\label{thm:g12:special-relativity:emc2}
\index{质量–能量等价}
质量为 $m$ 的物体，仅凭其质量便具有如下能量：
\[ E = m c^2 . \]
质量是能量的一种集中形式；反过来，能量也具有质量。
\end{theorem}
\admitted

\begin{remark}[你在哪里见过它]\label{rem:g12:special-relativity:emc2seen}
这正是核能背后的账目（见\cref{ch:g12:nuclear-energy}）：
反应的质量亏损乘以 $c^2$ 即为释放的能量。一
克质量对应 \qty{9.0e13}{J}——相当于一座大型电站一整天的发电量。
严格的推导留给大学分册；
反应堆与恒星不论如何都依靠它运转。
\end{remark}

\begin{remark}[牛顿并没有错]\label{rem:g12:special-relativity:newton}
取 $v \ll c$，则 $\gamma \to 1$，时间相同，长度不再收缩，
同时性成为绝对的——前三年所学的力学
完好无损地重现为低速极限。相对论并未推翻
牛顿，而是划定了他的理论疆域。物理学就是这样成长的：通过
不断完善理论，每一个新理论都把旧理论作为其极限情形保留下来。
\end{remark}

\section{练习}

\begin{exercise}[$\star$]\label{exo:g12:special-relativity:1}
求 $\gamma$ 的值（保留三位有效数字），取 $v = 0.1c$、$0.6c$、$0.8c$、$0.995c$。
\end{exercise}

\begin{exercise}[$\star$]\label{exo:g12:special-relativity:2}
一艘飞船以 $0.6c$ 巡航，其上的节拍器按飞船上的
时钟每 \qty{10.0}{s} 滴答一次。地面测得的时间间隔是多少？哪一个是固有时？为什么？
\end{exercise}

\begin{exercise}[$\star$]\label{exo:g12:special-relativity:3}
速率为多少（以 $c$ 的分数表示）时 $\gamma = 2$？又 $\gamma = 10$ 时呢？
\end{exercise}

\begin{exercise}[$\star$]\label{exo:g12:special-relativity:4}
一辆汽车以 \qty{130}{km/h} 行驶。求 $v/c$，再求 $\gamma - 1$
（见\cref{rem:g12:special-relativity:everyday}）；
汽车需行驶多久，其钟才落后 \qty{1}{s}？
\end{exercise}

\begin{exercise}[$\star$]\label{exo:g12:special-relativity:5}
在爱因斯坦的火车实验中，请引用所用的假设，解释
为何乘客断定\emph{前端}的闪击先发生，以及为何
两位观察者都没有错。
\end{exercise}

\begin{exercise}[$\star\star$]\label{exo:g12:special-relativity:6}
带电 π 介子的寿命为 $\Delta t_0 = \qty{2.6e-8}{s}$。一束
π 介子以 $0.99c$ 的速率离开加速器。求 $\gamma$、实验室中的
寿命、经典射程 $v\,\Delta t_0$ 以及实际平均射程。
\end{exercise}

\begin{exercise}[$\star\star$]\label{exo:g12:special-relativity:7}
由 $\gamma$ 曲线读出使 $\gamma = 1.5$ 和
$\gamma = 3$ 的速率（用计算检验）。$v = c$ 处的竖直渐近线
对有质量物体的加速说明了什么？
\end{exercise}

\begin{exercise}[$\star\star$]\label{exo:g12:special-relativity:8}
一个光钟的两面镜子相距 $L = \qty{1.5}{m}$。求其静止时的
滴答间隔，再求它以 $0.8c$ 巡航时地面所测得的滴答间隔。
\end{exercise}

\begin{exercise}[$\star\star$]\label{exo:g12:special-relativity:9}
谁测得以下事件之间的固有时：（a）腕表的两次滴答——是
佩戴者，还是路人？（b）μ子的诞生与衰变——是 μ子，还是
实验室？（c）离开地球与抵达火星——是驾驶员，还是任务控制中心？
\end{exercise}

\begin{exercise}[$\star\star$]\label{exo:g12:special-relativity:10}
空间站以 \qty{7.7}{km/s} 的速率绕行。求
$\gamma - 1 \approx v^2/2c^2$，再求六个月（\qty{1.58e7}{s}）后，
宇航员的钟比地面时钟落后多少。
\end{exercise}

\begin{exercise}[$\star\star$]\label{exo:g12:special-relativity:11}
求 \qty{1.0}{g} 物质所含的能量 $E = mc^2$，并将其
与一座 \qty{1.0}{GW} 电站一日的发电量相比，并与
一辆以 \qty{130}{km/h} 行驶的汽车所具有的 \qty{8.5e5}{J} 动能相比。
\end{exercise}

\begin{exercise}[$\star\star\star$]\label{exo:g12:special-relativity:12}
由 $\gamma = 1/\sqrt{1 - v^2/c^2}$ 证明 $v = c\sqrt{1 - 1/\gamma^2}$；
由此求 $\gamma = 30$ 的 μ子的速率（以
$c$ 为单位，保留五位有效数字）。
\end{exercise}

\begin{exercise}[$\star\star\star$]\label{exo:g12:special-relativity:13}
从 μ子自身的参考系出发复述其旅程：在 $\gamma = 30$ 时，大气层有多厚，
它以 $0.99944c$ 的速率掠过大气层需要多长时间，这段时间
是否在一个固有寿命之内？
\end{exercise}

\begin{exercise}[$\star\star\star$]\label{exo:g12:special-relativity:14}
一台直线加速器把电子加速到 $\gamma = \num{1.0e5}$，电子沿一根
长 \qty{3.2}{km} 的管道运动。电子在实验室中运动这段旅程需要多长时间？在
电子的参考系中需要多长时间，此时管道的“长度”又是多少？
\end{exercise}

\begin{exercise}[$\star\star\star$]\label{exo:g12:special-relativity:15}
比邻星位于 $4.2$ 光年之外（光年是光在一年内
所走过的距离）。一个探测器以 $0.9c$ 的速率飞往那里，
按地球上的日历和探测器上的时钟，旅程各需多长时间？
\end{exercise}

\section{问题：μ子与双胞胎}

\begin{problem}[{周末问题——μ子的旅程与宇航员的双胞胎：
一个比自己的“时钟”活得更久的粒子，重建的光钟，一位
比双胞胎年轻的旅行者，以及让 GPS 保持准确的核算}]
\label{pb:g12:special-relativity:1}
数据：μ子的固有寿命 $\Delta t_0 = \qty{2.2}{\micro s}$；μ子在
约 \qty{15}{km} 的高空以接近 $c$ 的速率产生；海平面的通量约为
每 \unit{cm^2} 每分钟约一个 μ子；$c = \qty{3.00e8}{m/s}$。

\medskip
\noindent\textbf{第一部分——μ子不可能的抵达。}
\begin{enumerate}
  \item 不考虑相对论，一个 μ子在一个寿命内能飞行多远？
  \item 那么，\qty{15}{km} 的下降需要经历多少个寿命？
  \item 经过 $n$ 个寿命后，存活的比例为 $\mathrm{e}^{-n}$：
        估算这一比例，并将其与实测通量相比较，用一句话对
        经典物理作出判断。
  \item 这些 μ子实际上以 $\gamma \approx 30$ 抵达。计算它们
        在地面测得的寿命，以及由此得到的射程。
  \item 利用 $v = c\sqrt{1 - 1/\gamma^2}$ 计算它们的速率，
        以 $c$ 为单位，保留五位有效数字。
\end{enumerate}

\medskip
\noindent\textbf{第二部分——重建光钟。}
一台光钟（$L = \qty{1.5}{m}$）装在以 $v = 0.6c$ 运动的飞船上。
\begin{enumerate}[resume]
  \item 计算船上读数所示的一次“滴答”时间 $\Delta t_0$。
  \item 从地面看，一次滴答持续 $\Delta t$：给出每半个滴答内的水平位移，
        光子的竖直爬升高度，以及——并指出决定其值的假设——
        斜向半程的长度。
  \item 对半程三角形应用毕达哥拉斯定理，解出
        $\Delta t$，并证明 $\Delta t = \gamma\,\Delta t_0$。
  \item 数值计算 $\gamma$ 与 $\Delta t$。
  \item 哪一位观测者测得两次滴答之间的固有时？
        请根据定义说明理由。
\end{enumerate}

\medskip
\noindent\textbf{第三部分——旅行者与她的双胞胎。}
一名宇航员离开留在地球上的双胞胎，以 $0.8c$ 的速率做一次往返旅行；
地球上的时钟测得整个旅程历时 \qty{10.0}{years}。
\begin{enumerate}[resume]
  \item 计算 $\gamma$，取 $0.8c$ 时的值（先给出精确分数，再给出小数）。
  \item 宇航员的时钟走过了多长时间？
  \item 折返点距离出发地有多远（以地球参考系度量，单位为光年）？
  \item 重逢时两人的年龄差是多少？哪一位更年轻？
  \item “可是运动是相对的——宇航员完全可以声称是地球在
        做这段旅行！”用一句话驳斥这一反对意见。
\end{enumerate}

\medskip
\noindent\textbf{第四部分——GPS 的账单。}
一颗 GPS 卫星沿半径为 $r = \qty{2.66e7}{m}$ 的轨道运行，周期为 $T = \qty{12}{h}$。
\begin{enumerate}[resume]
  \item 计算其轨道速率 $v = 2\pi r/T$。
  \item 计算 $\gamma - 1 \approx v^2/2c^2$。
  \item 以微秒为单位，推算卫星时钟每天相对地面时钟落后
        了多少（\qty{86400}{s}）。
  \item 若仅由这一偏差而不加修正，位置误差为
        $c \times$（偏差）：计算一天后的误差。
  \item 高处引力较弱，使时钟\emph{走快}
        \qty{45}{\micro s} 每天（广义相对论）：计算净偏差，
        以及以千米为单位的每日位置误差——这一数字迫使工程师将每台 GPS 时钟
        重新调校。
\end{enumerate}
\end{problem}

\bigskip
\noindent
这里是中学物理的终点，而不是物理学的终点。在大学分册中，
力学将借助微积分重新回归；电磁学
会显露为同一个结构，其中光与相对论本身，
一直隐藏其中；上一章初见端倪的量子世界，将成为
关于物质的成熟理论。每一个理论都将像相对论那样，
必须把你现在所认识的内容还原为一个极限情形。请带着它继续前进：这是
不会改变的那一部分。
